\documentclass[11pt,a4paper]{article}
% Shared preamble for RuPhO English problem statements (match T1 2026 style).
\usepackage[margin=1in]{geometry}
\usepackage{amsmath,amssymb}
\usepackage{graphicx}
\usepackage{float}
\usepackage{enumitem}
\usepackage{microtype}
\usepackage{caption}
\usepackage{hyperref}

\captionsetup{font=small,labelfont=bf,skip=6pt}

\setlist[itemize]{leftmargin=1.2cm}
\setlist[enumerate]{leftmargin=1.2cm}

\newcommand{\problembox}[3]{%
  \par\addvspace{0.42em}%
  \noindent
  \setlength{\fboxsep}{7.5pt}%
  \setlength{\fboxrule}{0.95pt}%
  \fbox{%
    \begin{minipage}{\dimexpr\linewidth-2\fboxsep-2\fboxrule\relax}
      \textbf{#1}\hfill\textbf{#3}\par\smallskip
      #2
    \end{minipage}%
  }%
  \par\addvspace{0.32em}%
}

\title{\textbf{Spaceship in the atmosphere}}
\author{\normalsize X25 (2025) — English translation}
\date{}
\begin{document}
\maketitle
Usually, when studying the dynamics of a body in a gravitational field, it is assumed that there are no frictional forces and the law of conservation of energy can be used. From this assumption one can obtain Kepler's laws, which describe with good accuracy both the behavior of planets and the movement of artificial satellites. However, at the initial and final stages of the spacecraft’s movement, it interacts with the planet’s atmosphere (if there is one), so air resistance forces significantly affect its trajectory. In this problem, we explore several simple models that allow us to evaluate the results of such influence. Further, for definiteness, we assume that the spacecraft is moving near the Earth.
A spacecraft moving in the atmosphere is subject to a gravitational force directed toward the center of the Earth, an air resistance force $D$ directed against the speed, and a lift force $L$ perpendicular to the speed. These forces are given by
\[
D = \frac{1}{2}\rho v^2 S C_D,\qquad
L = \frac{1}{2}\rho v^2 S C_L,
\]
where $\rho$ is the ambient air density, $S$ is the characteristic cross-sectional area of the spacecraft, $v$ is its speed relative to air, and $C_D$, $C_L$ are dimensionless coefficients. In general, these coefficients may depend on speed, but in this problem they are treated as constants. Assume that motion occurs in one fixed plane passing through the center of the Earth.

Use the atmospheric model
\[
\rho = \rho_0 e^{-\beta h},
\]
where $h$ is altitude above the Earth surface, $\rho_0$ is air density at the surface, and $\beta$ is a constant. Air can be considered stationary (ignore wind).

Throughout the problem you may use:
\[
\rho_0 = 1.225~\mathrm{kg/m^3},\quad
\beta = 0.1378~\mathrm{km^{-1}},\quad
r_0 = 6354~\mathrm{km},\quad
g_0 = 9.81~\mathrm{m/s^2}.
\]
The following numerical data can be used throughout the problem
\begin{figure}[H]
    \centering
    \includegraphics[width=0.72\textwidth]{T1_EN_assets/fig1.jpg}
    \caption{Fig. 1.}
\end{figure}
At a given moment in time, the motion of a spacecraft is characterized by the magnitude of its speed $v$, the angle $\gamma$ that the speed makes with the horizontal direction at the point where the spacecraft is located (it is considered positive if the speed is directed towards the surface of the Earth), as well as the height $h$ above the surface of the Earth. Let point $P$ be the projection of the ship’s position onto the Earth’s surface at some subsequent point in time. Then we define the horizontal displacement of the spacecraft $s$ as the distance between its initial and final projection, measured along the surface of the Earth (along a circular arc). You can also use the distance $r$ from the ship to the center of the Earth.
\section*{Part A. Equations of motion (1.2 points)}
\problembox{A1}{First, we obtain the exact equations that describe the dynamics of the spacecraft. Express the time derivatives of altitude $h$ and range $s$ in terms of $v$, $\gamma$, $r_0$, $r$.}{0.20}
\problembox{A2}{Find the derivatives of the velocity module $v$ and the angle $\gamma$. The response may also include the lift force $L$, the air resistance force $D$, the acceleration due to gravity at a given height $g$, the mass of the spacecraft $m$, $v$, $\gamma$, $r$.}{1.00}
\section*{Part B. Ballistic entry (2.8 points)}
In the general case, the resulting equations cannot be solved analytically. In order to get a qualitative idea of ​​the flight parameters, in this part we will consider the following simplified model. We will assume that there is no lift force ($L = 0$), and the angle $\gamma$ remains constant. Then we only need to investigate the dependence of the velocity modulus on height. In this case, we can assume that the force of air resistance is much greater than the component of gravity along the velocity, so gravity can also be ignored.
\problembox{B1}{Obtain an expression for the derivative $dv/dh$. Present the answer as
\[
\frac{dv}{dh}=A f(v,h),
\]
where $A$ is a constant that may depend on $m$, $S$, $C_D$, $C_L$, $\rho_0$, $\beta$, $\gamma$, and $f$ is some function of altitude and speed. Use constant $\mathbf{A}$ in later answers.}{0.40}
\problembox{B2}{Find the dependence of the speed of the spacecraft on altitude. At the initial altitude $h_0$ the speed was equal to $v_0$. Express your answer in terms of $A$, $\beta$, $h$, $v_0$, $h_0$. Also get an approximate answer, assuming that at the altitude $h_0$ the atmospheric density is negligible (in this case the answer should not include $h_0$). In the following, use an approximate expression throughout.}{0.80}
\problembox{B3}{Find the dependence of the acceleration $a$ (that is, the derivative of the velocity module) of the spacecraft on altitude. Express your answer in terms of $v_0$, $A$, $\beta$, $h$, $\gamma$.}{0.30}
\problembox{B4}{Find the height $h_c$ at which the acceleration modulus is maximum. Obtain a formula for the maximum value of acceleration modulus $a_{\text{max}}$ and velocity $v_c$ at height $h_c$. Express your answers in terms of $A$, $\beta$, $\gamma$, $v_0$.}{0.50}
The Vostok spacecraft was used for the first manned flight into space. For the reentry module, use:
\[
m = 2460~\mathrm{kg},\quad d = 2.3~\mathrm{m},\quad C_D = 2.0,\quad
\gamma = 3.2^\circ,\quad v_0 = 7.74~\mathrm{km/s},\quad h_0 = 122~\mathrm{km}.
\]
\problembox{B5}{For the given numerical data, find the speed (in km/s) and acceleration (in units of $g_0$) at heights $h_1 = 80~\mathrm{km}$, $h_2 = 60~\mathrm{km}$, $h_3 = 40~\mathrm{km}$.}{0.60}
\problembox{B6}{Find the maximum acceleration (in units of $g_0$) and the critical height $h_c$ at which it is achieved.}{0.20}
\section*{Part C. Sliding entry (2.2 points)}
In the previous part, we obtained quite large maximum acceleration values, which can lead to the destruction of the spacecraft. This acceleration can be significantly reduced by using a ship with significant lift. In order to describe the behavior of such a ship, the following approximations can be used: The angle $\gamma$ is small, $\sin \gamma \approx \gamma$, $\cos \gamma \approx 1$ The angle $\gamma$ changes very slowly during the movement, so that the contribution of the terms with $\dot{\gamma}$ in the equations of motion can be neglected. The acceleration of free fall is constant and equal to the acceleration on the Earth's surface $g_0$ The distance to the center of the Earth can be considered equal to the radius of the Earth $r \approx r_0$
In the previous part, we obtained quite large maximum acceleration values, which can lead to the destruction of the spacecraft. This acceleration can be significantly reduced by using a ship with significant lift. In order to describe the behavior of such a ship, the following approximations can be used:
\problembox{C1}{Find the first escape velocity $v_s$ – the speed of the spacecraft in a circular orbit, the radius of which is equal to the radius of the Earth. Express your answer in terms of $r_0$, $g_0$.}{0.10}
\problembox{C2}{Determine the speed at which the spacecraft must move at altitude $h$ for the motion described in the introduction to this part. Express your answer in terms of $v_s$, $\rho_0$, $\beta$, $C_L$, $S$, $m$, $r_0$, $h$.}{0.40}
\problembox{C3}{Determine the acceleration $a$ (derivative of the velocity modulus) created by the force of air resistance during such movement. Express your answer in terms of $v$, $C_L$, $C_D$, $r_0$, $g_0$.}{0.40}
\problembox{C4}{Let the initial speed of the spacecraft be $v_1< v_s$ and the final speed $v_2 < v_1$. Find the horizontal displacement of the ship $s$ during the movement. Consider that the entire change in the velocity module occurs due to the force of air resistance. Express your answer in terms of $r_0$, $C_L$, $C_D$, $v_s$, $v_1$, $v_2$.}{0.50}
\problembox{C5}{Let the following parameters be given for the spacecraft:
\[
m = 84\times10^3~\mathrm{kg},\quad S = 250~\mathrm{m^2},\quad C_D = 0.8,\quad C_L = 0.9.
\]
Find the speed (in km/s) and acceleration due to air resistance (in units of $g_0$) at altitudes $h_1 = 80~\mathrm{km}$, $h_2 = 60~\mathrm{km}$, $h_3 = 40~\mathrm{km}$.}{0.60}
\problembox{C6}{Under the conditions of the previous paragraph, calculate the range of horizontal movement of the spacecraft when descending from a height $h_0 = 90~\mathrm{km}$ to a height $h_f = 30~\mathrm{km}$.}{0.20}
\section*{Part D. Braking in the upper atmosphere (3.8 points)}
Let's consider the process of decreasing the orbit of a ship when moving in the upper atmosphere. The ship is acted upon by a drag force from the rarefied gas $\vec F=-\alpha \vec v$. For the considered motion in the upper layers of the atmosphere, the parameter $\alpha$ can be considered approximately constant. The force is very small, so we can assume that during one revolution the energy and angular momentum of the ship change very little.
Let at the moment of time $t$ a ship of mass $m$ move along an ellipse with semimajor axis $a$ and eccentricity $e$. Earth mass $M$.
\problembox{D1}{Write down expressions for energy $E$ and angular momentum $L$. Express your answer in terms of $G, m, a, e, M$. Note: It may be convenient to write the ZSE at the apocentric or periapsis point.}{0.40}
Note: It may be convenient to write the ZSE at the apocentric or periapsis point.
Let us begin to consider the effect of friction on the orbit of the ship.
\problembox{D2}{Write down the moment equation and expression for the power of the friction force $P$. Express your answer in terms of $\vec r, \vec v, \alpha$.}{0.20}
\problembox{D3}{Obtain an expression for the dependence of $\vec L$ on $t$ if at time $t=0$, $\vec L=\vec L_0$. Express your answer in terms of $\alpha$, $m$, $\vec L_0$.}{0.20}
To find the dependence of the ship's energy $E$ on time $t$, it is necessary to average the power of the friction force $P$ over the period of rotation of the ship. Since the energy and angular momentum change very little during one revolution, we will average the square of the velocity for orbital motion without taking into account the friction force. To do this, we use the Laplace–Runge–Lenz vector.
\problembox{D4}{Show that
\[
[\vec a,\vec L]=\beta\frac{d}{dt}\vec e_i.
\]
Here $\vec a$ is spacecraft acceleration and $\vec e_i$ is a suitable unit vector in polar coordinates centered at Earth. Find $\vec e_i$ and $\beta$. Note: the expression $\vec L=mr^2\vec\omega$ may be useful.}{0.30}
Note: It may be convenient to use the expression $\vec L=mr^2 \vec\omega$.
By integrating this expression over time, we can obtain
\[
[\vec v,\vec L]=\beta\vec e_i+\vec A \quad (1),
\]
where $\vec A$ is the Laplace-Runge-Lenz vector, constant for orbits in a gravitational field. Its magnitude is related to eccentricity by $|\vec A|=\beta e$.
\problembox{D5}{Multiplying expression (1) by $\vec r$ scalarly and using the results of point D4, obtain the dependence of the modulus of the radius vector $\vec r$ on $\varphi$ – the angle between the vectors $\vec A$ and $\vec e_r$. Express your answer in terms of $L, m, G, e, \varphi, M$.}{0.30}
\problembox{D6}{Using expression (1) and the results of point D4, obtain the dependence of the squared speed $v^2$ on $\varphi$. Express your answer in terms of $L, G, m, e, \varphi, M$.}{0.50}
\problembox{D7}{Obtain an expression for the time-average squared velocity $\langle v^2\rangle $ as an integral over the angle $\varphi$. Express your answer in terms of $  L,  m,  e,  M, G, a, \varphi, \tau$, the rocket’s orbital period. Note: It may be convenient to express the square of the average velocity from the averaged ZSE ($\langle T+U\rangle=\langle E\rangle$, where $T$ is the kinetic energy, $U$ is the potential energy).}{0.40}
Note: It may be convenient to express the square of the average velocity from the averaged ZSE ($\langle T+U\rangle=\langle E\rangle$, where $T$ is the kinetic energy, $U$ is the potential energy).
\problembox{D8}{Derive an expression for $ \langle\dot E\rangle $. Express your answer in terms of $G, m, a, e, M$.}{1.00}
\problembox{D9}{Using the results of points D1, D3, D8, obtain the dependences $a$ and $e$ on $t$, if at time $t=0$, $a=a_0$, $e=e_0$.}{0.40}
Let us consider the motion of the ISS in the upper atmosphere, which began in a circular orbit at altitude $h_0$.
\problembox{D10}{Find the time of descent of the orbit $T$ from height $h_0=408~\mathrm{km}$ to height $h=400~\mathrm{km}$, if $\alpha=7.17\cdot 10^{-5}~\mathrm{kg/s}$, $m=420~\mathrm{t}$, $r_0=6378~\mathrm{km}$. Express your answer in days.}{0.10}
\vspace{1em}
\noindent\textit{Source:} \url{https://pho.rs/p/4192}
\end{document}
